%% Full context-category and categorical beta-eta proofs.

\section{The cartesian closed category of simply typed contexts}

\begin{annexitem}{annex:p10.lambda.1}{thm:p10-stlc-equational-substitution}{Capture-avoiding simultaneous substitution}
\begin{proof}
All terms and contexts in a particular derivation are finite.
The infinite variable supply therefore permits bound variables
to be renamed, retaining their type annotations, outside all
names that the following constructions must avoid.

First prove weakening and consistent renaming for typing by
induction on the typing derivation.  A variable declaration
remains present; application uses the two induction hypotheses;
at abstraction, rename the bound variable if necessary, apply
the induction hypothesis to the extended context, and abstract
again.  The same induction on an equality derivation proves
weakening and renaming for equality.  Reflexivity, symmetry,
transitivity, and congruence are preserved.  A renamed or
weakened \(\beta\) equation is again such an equation, up to
bound renaming.  In an \(\eta\) equation choose its displayed
bound variable fresh also for the new declarations.

Define substitution recursively on terms.  For a variable
\(y_j\) in the source typing context, put \(y_j[\sigma]=\sigma_j\);
substitution commutes with application.  At an abstraction
\(\lambda z{:}\gamma.P\), make \(z\) fresh and use the lifted
substitution
\[
 \sigma^+:\Theta,z{:}\gamma\longrightarrow\Psi,z{:}\gamma,
 \qquad
 \sigma^+=(\sigma_1,\ldots,\sigma_n,z),
\]
where the old components have been weakened into
\(\Theta,z{:}\gamma\).  Set
\[
       (\lambda z{:}\gamma.P)[\sigma]
                =\lambda z{:}\gamma.P[\sigma^+].
\]
The definition does not depend, up to bound
\(\alpha\)-equivalence, on the fresh name: consistently
renaming it in the two recursively obtained bodies identifies
the results.  Induction on the term now proves typing.
The abstraction step applies the induction hypothesis to
\(P\) and \(\sigma^+\), then the abstraction typing rule.

The same recursion proves that an identity substitution acts
as the identity up to bound renaming.  For associativity,
freshen binders for the entire finite diagram of substitutions.
The variable case of
\[
                   (M[\sigma])[\tau]=_\alpha M[\sigma[\tau]]
\]
is the definition of the component tuple on the right, and
the application case follows componentwise.  In the abstraction
case both sides have the same fresh binder; their bodies agree
by induction and by
\(\sigma^+[\tau^+]=_\alpha(\sigma[\tau])^+\).
This proves the substitution identities.

Next induct on an equality derivation to substitute into it.
Equivalence rules and application congruence are immediate.
Abstraction congruence uses the lifted substitution.
For a \(\beta\) equation, its substituted left side contracts
to
\[
                 M[\sigma^+][N[\sigma]/x].
\]
Substitution associativity identifies this, up to bound
renaming, with \((M[N/x])[\sigma]\): both composite
substitutions send the old context variables to their
\(\sigma\)-components and \(x\) to \(N[\sigma]\).
For an \(\eta\) equation, choose \(x\) fresh for the context
and all components of \(\sigma\).  Its image is
\(\lambda x{:}\alpha.M[\sigma]x=M[\sigma]\), an instance
of the same typed schema.  Bound renaming is preserved by
capture avoidance.

Finally suppose \(\sigma_j=_{\beta\eta}\sigma'_j\) in their
common context for every component.  Induction on \(M\)
proves \(M[\sigma]=_{\beta\eta}M[\sigma']\).
The variable case is the corresponding hypothesis,
application uses congruence, and abstraction uses weakening
of the component equations, reflexivity for the fresh variable,
the induction hypothesis for the lifted substitutions, and
abstraction congruence.
\end{proof}
\end{annexitem}

\begin{annexitem}{annex:p10.lambda.2}{thm:p10-stlc-context-ccc}{Products and curried-list exponentials}
\begin{proof}
Substitution preserves typing, so composition produces an
arrow with the required target types.  If either input tuple
is replaced by a componentwise equal tuple,
Theorem~\ref{thm:p10-stlc-equational-substitution} gives the
same output class.  Consistent changes of the declared names
are bijective variable substitutions; composing with their
inverse restores the original presentation by the substitution
identities.  Thus intermediate choices of fresh names do not
affect composition.  The two identity laws and associativity
are the corresponding substitution identities, component by
component.  This establishes the category.

For any \(\Gamma\), there is exactly one tuple with no
components, so \(()\) is terminal.  In the concatenated
context \(\Gamma\cdot\Delta\), take the first and second
blocks of variables as the projections.  Pairing two arrows
with a common domain concatenates their tuples.
Substituting this tuple into a projection recovers its block.
Conversely, the two blocks determine every component of a
tuple into \(\Gamma\cdot\Delta\), proving the product
universal property.  These statements also apply when
either block is empty.

For the exponential construction, repeated application
types each \(f_jx_1\cdots x_m\) as \(\beta_j\).
Repeated abstraction types the proposed transpose component
as \(c_\Gamma(\beta_j)\).  Congruence makes transposition
well-defined on classes.  For an arrow \(K:\Xi\to\Delta^\Gamma\),
evaluation after \(K\times\mathrm{id}_\Gamma\) has components
\[
                         K_jx_1\cdots x_m.
\]
Consequently, evaluating the transpose of
\(T:\Xi\cdot\Gamma\to\Delta\) applies
\(\lambda x_1{:}\alpha_1\cdots\lambda x_m{:}\alpha_m.T_j\)
to the \(m\) corresponding variables.  Successive typed
\(\beta\) equations give \(T_j\), for every \(j\).

Conversely, transposing the displayed evaluation tuple gives
\[
          \lambda x_1{:}\alpha_1\cdots
                        \lambda x_m{:}\alpha_m.K_jx_1\cdots x_m.
\]
The \(x_i\)'s have been chosen outside the context of \(K_j\).
Apply \(\eta\) first to the innermost abstraction, then
successively to the remaining abstractions.  Each variable
is absent from the free variables of the term to which that
\(\eta\) step is applied, so the result is \(K_j\).
The two operations are inverse, which also proves uniqueness
of transposes.

If \(m=0\), the exponential object is \(\Delta\), evaluation
is its identity, and transposition changes no component.
If \(n=0\), the exponential object is \(()\), both hom-sets
in the exponential bijection are singletons, and all equations
are equations between empty tuples.  In particular the case
\(\Gamma=\Delta=()\) is included.

For \(r:\Xi'\to\Xi\), freshen the abstraction variables
outside its components.  Substitution under the \(m\)
abstractions then gives, componentwise,
\[
 \Lambda\bigl(T\circ(r\times\mathrm{id}_\Gamma)\bigr)
                  =\Lambda(T)\circ r.
\]
This is naturality in \(\Xi\).  Terminal object, products,
and these exponential bijections prove cartesian closure.
\end{proof}
\end{annexitem}

\begin{annexitem}{annex:p10.lambda.3}{thm:p10-stlc-ccc-completeness}{Categorical substitution and the separating syntactic model}
\begin{proof}
Interpret a named context by its ordered list of types.
Projections therefore give interpretations unchanged by
consistent renaming of the declared names.  Induction on terms,
using their annotated binders, proves bound
\(\alpha\)-invariance.  Typing of every interpretation follows
from the product and exponential domains and codomains in
Definition~\ref{def:p10-stlc-ccc-interpretation}.

We spell out the naturality of transposition used below.
For \(f:X\times A\to B\) and \(s:Y\to X\), the exponential
universal property gives
\[
        \Lambda(f\circ(s\times\mathrm{id}_A))
                             =\Lambda(f)\circ s.             \tag{*}
\]
Indeed, evaluation of the right side, after pairing it with
the \(A\)-projection, is \(f\circ(s\times\mathrm{id}_A)\),
by the evaluation equation for \(\Lambda(f)\) and the product
projection laws.  Uniqueness of transposes proves \((*)\).

First consider substitutions whose components are variables
of the corresponding types.  Their tuple arrows select the
indicated context projections.  Induction on the term proves
the substitution formula for these variable substitutions:
variables use the projection equations and application uses
pairing.  At an abstraction, the lifted selection retains the
fresh bound variable, so its tuple arrow is the old selection
arrow times the identity; apply the induction hypothesis to
the body and then \((*)\).
In particular, weakening a term to a context with one
additional final declaration interprets it by precomposition
with the projection to the old context.

We now prove the substitution formula by induction on \(M\).
For a variable it is exactly the defining projection equation
for the tuple arrow \(\llbracket\sigma\rrbracket\).
The application step uses the induction hypotheses and
\[
                 \langle f,g\rangle\circ s
                         =\langle f\circ s,g\circ s\rangle.
\]
For abstraction, use a fresh bound variable of type \(\gamma\).
The lifted substitution has interpretation
\(\llbracket\sigma\rrbracket\times
 \mathrm{id}_{\llbracket\gamma\rrbracket}\):
its old components are weakened interpretations and its last
component is the new projection.  The induction hypothesis
for the body and \((*)\) now give
\[
 \llbracket(\lambda x{:}\gamma.P)[\sigma]\rrbracket
  =\Lambda\bigl(\llbracket P\rrbracket\circ
       (\llbracket\sigma\rrbracket\times
        \mathrm{id}_{\llbracket\gamma\rrbracket})\bigr)
  =\llbracket\lambda x{:}\gamma.P\rrbracket
          \circ\llbracket\sigma\rrbracket.
\]
The empty target context has its unique terminal tuple arrow,
so this argument includes substitutions into closed terms.

For soundness of \(\beta\), let
\(f=\llbracket M\rrbracket:X\times A\to B\) and
\(n=\llbracket N\rrbracket:X\to A\).
The exponential evaluation equation gives
\[
           \operatorname{ev}\circ\langle\Lambda(f),n\rangle
                      = f\circ\langle\mathrm{id}_X,n\rangle.
\]
The right side is the interpretation of \(M[N/x]\) by the
substitution formula.  For \(\eta\), if \(m:X\to B^A\),
the transpose of
\(\operatorname{ev}\circ
 \langle m\circ\pi_X,\pi_A\rangle\) is \(m\), by the
uniqueness part of the exponential universal property.
This is exactly the interpretation of the typed \(\eta\)
equation, since the bound variable is fresh.
Equivalence rules, application congruence, and abstraction
congruence preserve equality of arrows.  Thus every typed
\(\beta\eta\) equation holds in every such interpretation.

For the converse take \(\mathcal C=\mathcal S_{\beta\eta}\)
with the precise choices of
Theorem~\ref{thm:p10-stlc-context-ccc}, and assign to each
base type \(\mathsf b\) the singleton context \((\mathsf b)\).
Induction on types gives
\[
                        \llbracket\alpha\rrbracket=(\alpha):
\]
the exponential of \((\beta)\) by \((\alpha)\) in this category
is the singleton
\((\alpha\mathbin{\to_{\mathrm{ty}}}\beta)\).
Starting from \(()\), recursive context products are
concatenations, so
\(\llbracket(\alpha_1,\ldots,\alpha_m)\rrbracket
   =(\alpha_1,\ldots,\alpha_m)\) literally.

Induction on a typed term shows that its interpretation in
this model is the arrow to the singleton target represented
by that very term.  Variables are the stated projections.
For application, pairing the two singleton arrows and
composing with the evaluation term \(fx\) produces \(MN\)
by simultaneous substitution.  For abstraction, the
transpose is the one-component tuple
\((\lambda x{:}\alpha.M)\) by construction.
Therefore equality of the interpretations of \(M\) and \(N\)
in this one CCC is precisely their componentwise typed
\(\beta\eta\) equality.  Equality in every CCC implies equality
in this CCC and proves the converse.  The argument uses no
normal-form theorem.
\end{proof}
\end{annexitem}
